\chapter{ANCA (Antineutrophil Cytoplasmic Antibodies)}

\section*{What Is This Test?}
ANCA testing detects antibodies directed mainly against proteinase 3 (PR3) and myeloperoxidase (MPO). It supports evaluation of ANCA-associated vasculitis (AAV), especially granulomatosis with polyangiitis (GPA) and microscopic polyangiitis (MPA), when the clinical presentation provides a reasonable pretest probability.

\section*{Alternative Names}
ANCA; anti-PR3; anti-MPO; c-ANCA; p-ANCA; vasculitis antibody testing.

\section*{Principle}
The 2017 international consensus recommends high-quality PR3- and MPO-specific immunoassays as the preferred initial screening method for suspected GPA or MPA. Indirect immunofluorescence (IIF) can be used when antigen-specific assays are negative despite high clinical suspicion, or to help evaluate low-positive or atypical results. IIF patterns are supportive but do not identify a disease by themselves \cite{bossuyt2017revised}.

\section*{History}
ANCA patterns and PR3/MPO specificity were developed through immunofluorescence and solid-phase immunoassays. Improved antigen-specific assays led to the revised 2017 international consensus, which moved away from a categorical IIF-first algorithm.

\section*{Why This Test Is Ordered}
\begin{itemize}
  \item Evaluation of rapidly progressive glomerulonephritis, pulmonary-renal syndrome, unexplained pulmonary hemorrhage, or compatible multisystem vasculitis.
  \item Supportive evaluation of GPA, MPA, or EGPA when clinical features are compatible.
  \item Evaluation of selected drug-induced vasculitis syndromes.
  \item Not for indiscriminate screening or diagnosis of vasculitis without a compatible clinical presentation.
\end{itemize}

\section*{Primary vs Secondary Test}
PR3/MPO immunoassay is a primary supportive test when clinical suspicion for AAV is appropriate. A positive result is not diagnostic by itself and does not automatically replace tissue diagnosis. Biopsy is pursued when needed to establish the diagnosis, define organ involvement, or exclude competing disease; treatment decisions in a life-threatening presentation should not be delayed solely to await a biopsy when specialist assessment supports immediate therapy.

\section*{How This Test Is Performed}
A venous blood sample is collected into a serum tube and analyzed by an antigen-specific immunoassay, with IIF or a second assay used selectively. Specimen requirements and cutoffs vary by laboratory.

\section*{What Preparation Is Needed}
Fasting is not required. Record immunosuppressive therapy and potentially causative drugs such as hydralazine, propylthiouracil, minocycline, or levamisole-adulterated cocaine exposure.

\section*{Contraindications and Precautions}
There are no absolute contraindications to venipuncture. A positive PR3 or MPO result contributes to the diagnostic workup but is not diagnostic alone. A negative result does not exclude AAV. IIF patterns can be positive in inflammatory bowel disease, autoimmune liver disease, infection, and drug-induced autoimmunity. Serial ANCA titers should not be used alone to guide treatment or define relapse.

\section*{Risks}
Minimal venipuncture risks include brief pain, bruising, hematoma, and vasovagal symptoms.

\section*{What Happens After the Test}
Interpret the result with urinalysis, creatinine/eGFR, inflammatory markers, imaging, symptoms, and examination. PR3 positivity supports GPA and MPO positivity supports MPA or EGPA in the appropriate clinical setting, but neither is sufficient to establish the diagnosis. Kidney, skin, lung, nerve, or other tissue biopsy may be required. Urgent specialist assessment is needed for suspected pulmonary-renal syndrome or rapidly progressive organ injury.

\section*{Understanding Results}
\subsection*{Below Normal (Negative)}
A negative PR3/MPO assay does not exclude AAV, particularly EGPA or limited disease. If clinical suspicion remains high, consult the laboratory and specialist team about IIF, another antigen-specific assay, and tissue diagnosis.

\subsection*{Above Normal (Positive)}
A positive PR3 or MPO result supports AAV only when the clinical presentation is compatible. Low-positive or atypical results require confirmation and consideration of nonvasculitic causes. Do not diagnose or treat solely from an ANCA result.

\section*{Normal Results}
No PR3 or MPO antibody detected by the reporting laboratory's assay. IIF, when performed, is reported as negative according to that laboratory's method and cutoff.

\section*{Follow-up Tests}
\begin{itemize}
  \item Urinalysis with microscopy, urine protein quantification, creatinine/eGFR, CBC, and inflammatory markers.
  \item Anti-GBM antibody and other serologies when evaluating pulmonary-renal syndrome.
  \item Chest imaging and organ-specific evaluation.
  \item Tissue biopsy when needed to establish diagnosis or exclude mimics.
  \item Clinical monitoring rather than ANCA titer alone for treatment response and relapse.
\end{itemize}

\section*{References}
\bibliographystyle{unsrtnat}
\bibliography{references}
\clearpage
\section*{Regulatory Status and Authorization Date}
Regulatory status applies to the specific assay, instrument, manufacturer, and intended use rather than to the test name alone. No single approval date can be assigned to this test category. For a commercial assay, verify the current product in the relevant regulator's database: the U.S. Food and Drug Administration (FDA) clearance, approval, or authorization; India's Central Drugs Standard Control Organisation (CDSCO) licence or permission; the European Union In Vitro Diagnostic Regulation (EU IVDR) CE certificate issued through the applicable conformity-assessment route; or the United Kingdom Medicines and Healthcare products Regulatory Agency (MHRA) registration and UKCA/CE status. Laboratory-developed tests may not have an individual FDA, CDSCO, EU IVDR, or MHRA approval date.
\clearpage
